% Vocabulary-Set Interference in Open-Vocabulary Object Detection. Draft. Every number must come from ovd/RESULTS.md;
% missing numbers are [TODO]. Build: pdflatex paper && pdflatex paper (inside ovd/paper).
\documentclass[10pt,twocolumn,letterpaper]{article}
\usepackage[review]{cvpr}
\usepackage{graphicx,booktabs,amsmath,amssymb}
\usepackage[pagebackref,breaklinks,colorlinks,allcolors=blue]{hyperref}
\def\paperID{*****}
\def\confName{ICCV}
\def\confYear{2027}
\newcommand{\todo}[1]{\textcolor{red}{[TODO: #1]}}
\title{Vocabulary-Set Interference in Open-Vocabulary Object Detection}
\author{Anonymous submission}
\begin{document}
\maketitle

\begin{abstract}
Open-vocabulary detectors are queried with a user-chosen list of class names. We ask a narrow question: for the
\emph{same} image and the \emph{same} target name, does the detector's answer change when \emph{other} names are added,
removed or reordered? We build a counterfactual protocol that holds the target phrase, image, preprocessing and
weights fixed and varies only the rest of the vocabulary (1 to 1{,}203 classes; random, semantically similar, visually
confusable and unrelated distractors; five orders), and apply it to Grounding DINO, YOLO-World and OWLv2 on COCO and
LVIS. \todo{main measured effect, with CI}. We trace the dependence to specific architectural steps: OWLv2's class
scores are invariant by construction and change only through top-1 label competition; YOLO-World's text guidance is a
maximum over the vocabulary, which makes it exactly order-invariant but set-dependent; Grounding DINO's text encoder is
masked per phrase, so its dependence arises in cross-modal fusion, query selection, label competition and, beyond 256
tokens, chunking. We then show training-free inference changes that are vocabulary-invariant by construction --
fixed scoring contexts, a fixed guide vocabulary and independent read-out -- and measure their accuracy and cost:
\todo{accuracy vs.\ the strongest baseline, latency}.
\end{abstract}

\section{Introduction}
Open-vocabulary detectors~\cite{gdino,yoloworld,owlv2} let a user name the classes to detect at test time. In practice,
the vocabulary is an engineering choice: a deployment lists the classes it cares about, a benchmark lists all of its
classes, and long vocabularies are split into chunks to fit text-encoder limits. Implicitly, users assume that the
answer for ``dog'' does not depend on whether ``cat'' or ``toaster'' is also in the list. Recent work shows that
whole-vocabulary choices matter -- pruning the vocabulary per image helps~\cite{vocada}, and Grounding DINO collapses on
aerial imagery when given all 80 classes rather than the few present~\cite{aerial} -- but these comparisons change the
vocabulary and the information about which classes are present at the same time, and do not say \emph{where} the
dependence arises.

We isolate the effect with a counterfactual protocol: the image, the target phrase, the preprocessing and the weights
are fixed, and only the other class names change. This lets us (i) measure interference per target instance --
score changes, detection flips at the deployment threshold, box changes, absent-class false alarms and AP; (ii)
attribute it to architectural steps, in some cases exactly; and (iii) remove it with training-free changes that are
invariant by construction, whose only open questions are accuracy and cost.

\paragraph{Contributions.} \todo{finalise after results}
(1) A fixed-target counterfactual protocol and benchmark for vocabulary interference on COCO and LVIS.
(2) A per-architecture mechanism analysis separating network-level interference from label competition, thresholds,
NMS and chunking.
(3) Training-free, vocabulary-invariant inference for each detector family, with accuracy and latency against original,
single-class and vocabulary-pruning baselines.

\section{Related work}
\paragraph{Open-vocabulary detection.} Grounding DINO~\cite{gdino} fuses image and text early and grounds phrases;
YOLO-World~\cite{yoloworld} injects text embeddings into a YOLO neck for real-time use; OWL-ViT/OWLv2~\cite{owlvit,owlv2}
score region embeddings against each query independently. Grounding DINO introduced per-phrase attention masks to avoid
``cross-talk'' between names~\cite{gdino}; we measure what cross-talk remains.
\paragraph{Vocabulary choice.} VocAda~\cite{vocada} prunes the user vocabulary per image with a captioner and a selector
(Detic, VLDet, CoDet; COCO, Objects365). An aerial evaluation~\cite{aerial} reports a large gap between all-class and
present-class prompting for Grounding DINO. Chunked inference over long vocabularies is standard and known to change
predictions~\cite{mmgdino,slground}. None of these hold the target fixed or attribute the effect to architectural steps.
\paragraph{Score reliability.} Open-vocabulary scores depend on object scale and on the wording of the target's own
query~\cite{scalesem}; confident false positives are common on background images~\cite{tractordogs}. We study the
complementary factor: the \emph{other} queries.

\section{Protocol}
\paragraph{Data.} COCO val2017~\cite{coco} and LVIS v1.0 val~\cite{lvis}, public annotations only. Images are ranked by an
MD5 hash of the file name: the first 500 COCO images are a development split (all design choices); the next 1{,}000
COCO images and 600 LVIS images (excluding the development images) are held out. Targets are every present COCO class
(LVIS: up to three) plus two absent classes (LVIS: verified absent).
\paragraph{Conditions.} For target $c$: $\{c\}$; $c$ plus $k{-}1$ random classes ($k\in\{5,20\}$ on COCO;
$\{5,20,80,200,500\}$ on LVIS); $c$ plus its nearest classes by CLIP text embedding (similar), plus the classes CLIP ranks
highest for the target's own crop (visually confusable), plus random classes from the least similar half (unrelated);
and the full list in five orders. The target string is identical in all conditions; seeds derive from hashes.
\paragraph{Measurements.} Raw target score of the best-matching box (IoU$\ge$0.5) per instance; deployed output
(each box keeps its top class and must exceed the model's threshold: 0.35 / 0.25 / 0.10 for Grounding DINO / YOLO-World
/ OWLv2); detection flips; best-box IoU change; false alarms on absent targets; AP50 per class over the target's pairs;
latency. Confidence intervals: bootstrap over images.

\section{Where does the vocabulary enter?}\label{sec:mech}
\paragraph{OWLv2.} The logit of box $b$ for query $q$ is a function of $b$'s embedding and $q$'s text embedding only,
so raw scores are exactly invariant to the other queries (we confirm numerically: \todo{max difference}). Any change
in deployed output comes from top-1 label assignment.
\paragraph{YOLO-World.} Text enters twice: as the classifier (a per-class contrastive head) and as a \emph{guide} in the
neck, where max-sigmoid attention weights each location by $\sigma(\max_{t\in V}\langle e_x, g(t)\rangle)$. The maximum
makes the features exactly order-invariant and monotone in the vocabulary set: adding a name can only raise the
attention at each location. Fixing the guide to a constant list makes every class score independent of $V$
(numerically, $<10^{-6}$ difference).
\paragraph{Grounding DINO.} Phrases are separated by attention masks and position ids restart per phrase in the text
encoder, so the target's text features do not depend on other names there. Other names enter through (a) image--text
fusion in the feature enhancer and decoder, (b) language-guided selection of the 900 decoder queries (maximum over all
tokens), (c) top-1 label competition and (d) chunking when the prompt exceeds 256 tokens, where order decides which
names share the target's chunk. We probe (b) by restricting query selection to the target's tokens.
\todo{results of the probe and of each factor}

\section{Vocabulary-invariant inference}
\textbf{Fixed context} (Grounding DINO): every class is always scored inside the same chunk of a fixed reference list in
canonical order; the deployment vocabulary only selects which classes are reported (COCO: one 195-token prompt).
\textbf{Fixed guide} (YOLO-World): the neck is guided by a fixed reference list, the head uses the deployment vocabulary.
\textbf{Independent read-out} (all): no top-1 label competition. These are invariant by construction, so we evaluate
accuracy and cost against original inference with the full list, single-class inference and CLIP-based vocabulary
pruning in the style of VocAda. \todo{table}

\section{Results}
\todo{pilot (development split), held-out COCO, LVIS, mechanism ablations, corrections, latency, error analysis}

\section{Limitations}
Base/small checkpoints only (one 24\,GB GPU); two datasets; the visually confusable family is defined with CLIP; the
pruning baseline omits VocAda's captioner. \todo{add result-specific limitations}

{\small
\begin{thebibliography}{99}
\bibitem{gdino} S. Liu, Z. Zeng, T. Ren, et al. Grounding DINO: Marrying DINO with grounded pre-training for open-set object detection. In \emph{ECCV}, 2024.
\bibitem{yoloworld} T. Cheng, L. Song, Y. Ge, W. Liu, X. Wang, Y. Shan. YOLO-World: Real-time open-vocabulary object detection. In \emph{CVPR}, 2024.
\bibitem{owlvit} M. Minderer, A. Gritsenko, A. Stone, et al. Simple open-vocabulary object detection with vision transformers. In \emph{ECCV}, 2022.
\bibitem{owlv2} M. Minderer, A. Gritsenko, N. Houlsby. Scaling open-vocabulary object detection. In \emph{NeurIPS}, 2023.
\bibitem{vocada} M. Liu, T. L. Hayes, M. Mancini, E. Ricci, R. Volpi, G. Csurka. Test-time vocabulary adaptation for language-driven object detection. In \emph{ICIP}, 2025. arXiv:2506.00333.
\bibitem{aerial} C. Tsourveloudis. Do open-vocabulary detectors transfer to aerial imagery? A comparative evaluation. arXiv:2601.22164, 2026.
\bibitem{scalesem} Y. T. Soon, J.-W. Hsieh. Confidence scores in open-vocabulary detection are a biased mixture of scale and semantics. arXiv:2607.10993, 2026.
\bibitem{slground} SL-Ground: Sentence-level grounding for fine-grained open-vocabulary object detection. arXiv:2610.24814, 2026. [VERIFY authors]
\bibitem{mmgdino} X. Zhao, Y. Chen, S. Xu, et al. An open and comprehensive pipeline for unified object grounding and detection (MM-Grounding-DINO). arXiv:2401.02361, 2024. [VERIFY]
\bibitem{tractordogs} Fantastic tractor-dogs and how not to find them. In \emph{ICLR}, 2026. [VERIFY authors and title]
\bibitem{coco} T.-Y. Lin, M. Maire, S. Belongie, et al. Microsoft COCO: Common objects in context. In \emph{ECCV}, 2014.
\bibitem{lvis} A. Gupta, P. Doll\'ar, R. Girshick. LVIS: A dataset for large vocabulary instance segmentation. In \emph{CVPR}, 2019.
\end{thebibliography}
}
\end{document}
